\chapter{Acotado y Anotaciones}

\section{Normas de Acotado y Anotaciones Técnicas}

El acotado (dimensionamiento) es el proceso crítico mediante el cual
se indican en el dibujo técnico las dimensiones reales de un objeto,
para que este pueda ser fabricado sin ambigüedades. Las cotas deben ser
precisas, legibles y seguir estrictas normas IRAM que regulan su
representación.

\subsection{Elementos de una Cota}

Toda cota consta de tres elementos fundamentales:
\begin{itemize}
	\item \textbf{Línea de cota:} Es la línea fina sobre la cual se
	      coloca el valor numérico.
	\item \textbf{Líneas auxiliares (o de referencia):} Son líneas
	      finas que se extienden desde el objeto hacia la línea de
	      cota.
	\item \textbf{Flechas o terminaciones:} Indican los límites de la
	      cota.
	\item \textbf{Valor de la cota:} Es la medida real del objeto,
	      que se expresa siempre en la misma unidad (generalmente mm).
\end{itemize}

\subsection{Reglas Fundamentales del Acotado}

\begin{enumerate}
	\item \textbf{Claridad:} Las cotas no deben entorpecer la
	      lectura del dibujo.
	\item \textbf{No redundancia:} Una cota solo debe aparecer una vez
	      en el plano.
	\item \textbf{Funcionalidad:} Se deben acotar las dimensiones
	      que son importantes para el funcionamiento o el montaje de
	      la pieza.
	\item \textbf{Lectura:} Los valores de las cotas deben poder
	      leerse desde abajo o desde la derecha del plano.
\end{enumerate}

\subsection{Anotaciones y Tolerancias}

Las anotaciones son textos complementarios que proporcionan información
adicional indispensable: materiales, tratamientos térmicos, acabados
superficiales, notas sobre procesos de fabricación, etc.

Por otro lado, la \textbf{tolerancia dimensional} es el margen de error
permitido en una dimensión. Ninguna pieza se fabrica con exactitud
matemática; la tolerancia garantiza que la pieza sea funcional dentro
de un rango aceptable.

\section{Cuestionario}
\begin{enumerate}
	\item Define cota, línea de cota y línea auxiliar.
	\item ¿Qué es la escala de cota y cuándo se ajusta?
	\item ¿Cómo deben ser las flechas en un acotado normalizado?
	\item Menciona las reglas fundamentales para evitar el acotado
	      redundante.
	\item ¿Qué es una cota funcional vs. una cota no funcional?
	\item ¿Cómo influye el acabado superficial en la pieza y cómo se
	      indica?
	\item ¿Qué es la tolerancia dimensional y por qué es necesaria?
	\item ¿Cómo se acotan correctamente arcos y círculos?
\end{enumerate}

\section{Parte Práctica}
\begin{enumerate}
	\item Configurar el estilo de cota en el software: fuente, altura de
	      texto, tipo de flecha y precisión decimal.
	\item Acotar un objeto simple utilizando cotas lineales horizontales
	      y verticales.
	\item Acotar un diámetro y un radio usando los símbolos adecuados.
	\item Crear una cota alineada para una cara inclinada de una pieza.
	\item Añadir un texto de nota en un plano (ej: "Tratamiento térmico:
	      templado").
	\item Acotar una pieza compleja utilizando al menos tres vistas
	      diferentes.
	\item Usar cotas continuas (en cadena) y paralelas (por base) para
	      una misma pieza.
	\item Crear una tolerancia dimensional para un eje (ej: 20 +/- 0.05
	      mm).
	\item Añadir un símbolo normalizado de acabado superficial en una cara.
	\item Realizar un acotado completo para un plano de fabricación
	      mecánica.
\end{enumerate}
